\section{Discussion, limitations, and threats to validity}
\label{sec:limitations}

\subsection{Scientific scope}

The strongest trainable evidence concerns small decoder models on controlled
retrieval and byte-level language modeling. It does not directly cover
pretrained open-weight language models, natural images, open RLVR, or deployed
agent populations. The cross-domain dossiers are model-validation cases and are
not coequal replications. A seven-coordinate taxonomy improves the precision of
these boundaries but cannot create missing external evidence.

The framework also does not assume that every system has an equilibrium Gibbs
measure or a scalar thermodynamic limit. Optimizer trajectories, diffusion
paths, feedback loops, and changing populations are non-equilibrium objects.
The appropriate result may be a crossover, dynamical regime, path dependence,
or censoring rather than a phase transition.

\subsection{Reference-grid limitations}

The GPU reference has only three widths, three nominal controls, five seeds,
and 24 update steps. One nominal control pair collapses to the same effective
data count, and all conditions reach the end of the short training horizon
before the registered time-to-order criterion. These features limit both
control resolution and dynamical interpretation. They are now shown directly
rather than hidden by a smooth line.

Outer-seed intervals quantify initialization/data/minibatch variability under
the registered generator but do not establish independent implementation
replication. With only five seeds, susceptibility intervals are wide and power
to detect small OOD penalties is limited. The negative phase result applies to
the tested grid; it is not a proof of absence at larger width or longer time.

\subsection{Exploratory multiplicity and confirmation}

The broad tensor evaluates many families and observables. Its five-seed
intervals are descriptive screens, not multiplicity-corrected hypothesis tests.
The twelve-seed suite protects selected numerical contrasts with fresh seeds,
but its hypotheses were chosen after inspecting the screen. It is a second-stage
confirmation, not an independent external replication.

Optimizer contrasts use frozen learning-rate slices but do not exhaustively
match update norm, noise scale, training compute, or tuning budget. Corpus bits
per byte are tokenizer independent, yet corpora differ in intrinsic
predictability and preprocessing. Neither panel should be read as a universal
ranking.

\subsection{Construct validity}

Performance-derived order remains a task order unless a semantic probe
connects it to a latent or functional property. Attention entropy is not
strategy entropy; spectral effective rank is not macrostate multiplicity; and
temporal checkpoint variance is not outer-disorder susceptibility. The
observable dictionary reduces these errors, but any misspecified latent probe or
null can still compromise construct validity.

Finite-size forms are empirical candidates outside solvable anchors. Width,
depth, data, context, trajectory resolution, horizon, and population need not
share an exponent. Data collapse can be overfit when too many exponents and
windows are selected on the same sizes; the atlas therefore requires a
prospective largest-size check before using phase language.

\subsection{Broader use and non-goals}

The atlas is intended to make claims easier to compare, reproduce, and falsify.
It should not be used to lend physical prestige to a benchmark curve or to rank
systems with a single evidence score. Theory, fidelity, evidence, and outcome
remain vectors because collapsing them would erase the very distinctions the
framework is designed to preserve.

\FloatBarrier
